\documentclass[aspectratio=169,10pt]{beamer}
% ═══════════════════════════════════════════════════════════════════════════════
% SHARED PREAMBLE — Foundation Models Course (HIT)
% To be \input by each lecture file AFTER \documentclass
% ═══════════════════════════════════════════════════════════════════════════════

% ─────────────────────────────────────────────────────────────────────────────
% BEAMER THEME & BASIC SETUP
% ─────────────────────────────────────────────────────────────────────────────
\usetheme{Madrid}
\usecolortheme{default}

% Remove navigation symbols
\beamertemplatenavigationsymbolsempty

% ─────────────────────────────────────────────────────────────────────────────
% COLORS — HIT Branding
% ─────────────────────────────────────────────────────────────────────────────
\definecolor{hitblue}{RGB}{0,51,102}        % Primary: HIT Navy
\definecolor{hitgold}{RGB}{192,148,0}       % Accent: HIT Gold
\definecolor{hitlight}{RGB}{230,240,250}    % Light background

% Override Beamer palette
\setbeamercolor{primary}{fg=white,bg=hitblue}
\setbeamercolor{secondary}{fg=hitblue,bg=hitlight}
\setbeamercolor{structure}{fg=hitblue}
\setbeamercolor{alerted text}{fg=hitgold}

% ─────────────────────────────────────────────────────────────────────────────
% PACKAGES
% ─────────────────────────────────────────────────────────────────────────────
\usepackage{amsmath}
\usepackage{amssymb}
\usepackage{mathtools}
\usepackage{booktabs}
\usepackage{graphicx}
\usepackage{hyperref}
\usepackage{tikz}
\usepackage{pgfplots}
\usepackage{tcolorbox}
\usepackage{listings}
\usepackage{xcolor}
\usepackage{algorithm}
\usepackage{algpseudocode}

% ─────────────────────────────────────────────────────────────────────────────
% TikZ LIBRARIES
% ─────────────────────────────────────────────────────────────────────────────
\usetikzlibrary{shapes,arrows,positioning,calc,chains,scopes}
\pgfplotsset{compat=1.17}

% ─────────────────────────────────────────────────────────────────────────────
% CODE LISTINGS
% ─────────────────────────────────────────────────────────────────────────────
\lstset{
  basicstyle=\ttfamily\small,
  breaklines=true,
  commentstyle=\color{gray},
  keywordstyle=\color{hitblue},
  stringstyle=\color{hitgold},
  showstringspaces=false,
  backgroundcolor=\color{hitlight},
  frame=single,
  rulecolor=\color{hitblue}
}

% ─────────────────────────────────────────────────────────────────────────────
% TCOLORBOX STYLES
% ─────────────────────────────────────────────────────────────────────────────
\tcbset{
  hitbox/.style={
    colback=hitlight,
    colframe=hitblue,
    fonttitle=\bfseries,
    boxrule=1pt
  },
  alertbox/.style={
    colback=yellow!10,
    colframe=hitgold,
    fonttitle=\bfseries,
    boxrule=1.5pt
  }
}

% ─────────────────────────────────────────────────────────────────────────────
% CUSTOM COMMANDS
% ─────────────────────────────────────────────────────────────────────────────

% Placeholder for figures
\newcommand{\placeholder}[3]{
  \begin{center}
    \fbox{\begin{minipage}[c][#2][c]{#1}
      \centering\small\color{gray}[Figure: #3]
    \end{minipage}}
  \end{center}
}

% Section divider frame
\newcommand{\sectionframe}[2]{
  \begin{frame}[plain]
    \begin{tikzpicture}[remember picture,overlay]
      \fill[hitblue] (current page.south west) rectangle (current page.north east);
      \node[anchor=center, white, font=\Large\bfseries] at (current page.center) {
        \begin{tabular}{c}
          #1 \\[0.3cm]
          {\small\color{hitgold}#2}
        \end{tabular}
      };
    \end{tikzpicture}
  \end{frame}
}

% ─────────────────────────────────────────────────────────────────────────────
% LOAD CUSTOM MACROS
% ─────────────────────────────────────────────────────────────────────────────
% ═══════════════════════════════════════════════════════════════════════════════
% SHARED MACROS — Mathematical Notation for Foundation Models Course
% ═══════════════════════════════════════════════════════════════════════════════

% ─────────────────────────────────────────────────────────────────────────────
% ACTIVATION & LAYER FUNCTIONS
% ─────────────────────────────────────────────────────────────────────────────
\newcommand{\softmax}{\mathrm{softmax}}
\newcommand{\sigmoid}{\mathrm{sigmoid}}
\newcommand{\relu}{\mathrm{ReLU}}
\newcommand{\gelu}{\mathrm{GELU}}
\newcommand{\LayerNorm}{\mathrm{LayerNorm}}
\newcommand{\FFN}{\mathrm{FFN}}
\newcommand{\MHA}{\mathrm{MHA}}
\newcommand{\Attn}{\mathrm{Attention}}

% ─────────────────────────────────────────────────────────────────────────────
% ATTENTION COMPONENTS
% ─────────────────────────────────────────────────────────────────────────────
\newcommand{\query}{Q}
\newcommand{\key}{K}
\newcommand{\val}{V}
\newcommand{\dmodel}{d_{\text{model}}}
\newcommand{\dk}{d_k}
\newcommand{\nheads}{h}

% Scaled dot-product attention formula
\newcommand{\SDPA}{
  \Attn(\query, \key, \val) = \softmax\left(\frac{\query \key^T}{\sqrt{\dk}}\right) \val
}

\newcommand{\CrossAttn}{\text{CrossAttention}}

% ─────────────────────────────────────────────────────────────────────────────
% PROBABILITY & EXPECTATION
% ─────────────────────────────────────────────────────────────────────────────
\newcommand{\E}{\mathbb{E}}
\newcommand{\Prob}{\mathbb{P}}
\newcommand{\KL}{\mathrm{KL}}
\newcommand{\ELBO}{\mathrm{ELBO}}
\newcommand{\given}{\mid}

% ─────────────────────────────────────────────────────────────────────────────
% NORMS & OPERATIONS
% ─────────────────────────────────────────────────────────────────────────────
\newcommand{\norm}[1]{\left\lVert #1 \right\rVert}
\newcommand{\abs}[1]{\left| #1 \right|}
\newcommand{\inner}[2]{\langle #1, #2 \rangle}
\newcommand{\T}{^{\top}}

% ─────────────────────────────────────────────────────────────────────────────
% VECTORS & MATRICES
% ─────────────────────────────────────────────────────────────────────────────
\newcommand{\bvec}[1]{\mathbf{#1}}
\newcommand{\mat}[1]{\mathbf{#1}}
\newcommand{\iid}{\stackrel{\text{i.i.d.}}{\sim}}
\newcommand{\defeq}{\coloneq}

% ─────────────────────────────────────────────────────────────────────────────
% LANGUAGE MODELING
% ─────────────────────────────────────────────────────────────────────────────
\newcommand{\vocab}{\mathcal{V}}
\newcommand{\context}{\text{context}}
\newcommand{\token}{t}
\newcommand{\seq}{x}
\newcommand{\ptheta}{p_\theta}
\newcommand{\pref}{p_{\text{ref}}}

% ─────────────────────────────────────────────────────────────────────────────
% RLHF & DPO
% ─────────────────────────────────────────────────────────────────────────────
\newcommand{\piref}{\pi_{\text{ref}}}
\newcommand{\pitheta}{\pi_\theta}
\newcommand{\yw}{y_w}
\newcommand{\yl}{y_l}
\newcommand{\DPOloss}{\mathcal{L}_{\mathrm{DPO}}}

% ─────────────────────────────────────────────────────────────────────────────
% RETRIEVAL & RAG
% ─────────────────────────────────────────────────────────────────────────────
\newcommand{\docs}{\mathcal{D}}
\newcommand{\qvec}{q}
\newcommand{\dvec}{d}

% ─────────────────────────────────────────────────────────────────────────────
% AGENTS & REASONING
% ─────────────────────────────────────────────────────────────────────────────
\newcommand{\obs}{o}
\newcommand{\act}{a}
\newcommand{\thought}{t}
\newcommand{\tools}{\mathcal{T}}
\newcommand{\history}{h}
\newcommand{\concat}{\oplus}

% ─────────────────────────────────────────────────────────────────────────────
% SETS & LOGIC
% ─────────────────────────────────────────────────────────────────────────────
\newcommand{\R}{\mathbb{R}}
\newcommand{\N}{\mathbb{N}}
\newcommand{\Z}{\mathbb{Z}}

\endinput


% ─────────────────────────────────────────────────────────────────────────────
% TITLE & AUTHOR (can be overridden in individual lectures)
% ─────────────────────────────────────────────────────────────────────────────
\author[D. Lederman]{Dror Lederman}
\institute{Holon Institute of Technology (HIT) \\ Data Science Department}
\date{\today}

% ─────────────────────────────────────────────────────────────────────────────
% FOOTER & HEADER
% ─────────────────────────────────────────────────────────────────────────────
\setbeamertemplate{footline}{
  \leavevmode%
  \hbox{%
    \begin{beamercolorbox}[wd=0.33\paperwidth,ht=2.25ex,dp=1ex,center]{author in head/foot}%
      \usebeamerfont{author in head/foot}\insertshortauthor
    \end{beamercolorbox}%
    \begin{beamercolorbox}[wd=0.34\paperwidth,ht=2.25ex,dp=1ex,center]{title in head/foot}%
      \usebeamerfont{title in head/foot}\insertshorttitle
    \end{beamercolorbox}%
    \begin{beamercolorbox}[wd=0.33\paperwidth,ht=2.25ex,dp=1ex,right]{frame number}%
      \usebeamerfont{frame number}\insertframenumber{} / \inserttotalframenumber\hspace*{2ex}
    \end{beamercolorbox}%
  }%
  \vskip0pt%
}

\endinput


\title{Week 13: Future Directions and Student Presentations}
\subtitle{Foundation Models and Agentic AI}
\author{Lecturer Name}
\institute{Holon Institute of Technology (HIT)}
\date{Semester, 2025}

\begin{document}

\begin{frame}
  \titlepage
\end{frame}

\begin{frame}{Session Overview}
  \begin{columns}[T]
    \column{0.50\textwidth}
      \textbf{Part 1 (60 min): Future Directions}
      \begin{itemize}
        \item Autonomous research agents
        \item AI scientists
        \item Agentic workflows in production
        \item Open research problems
        \item Societal implications
      \end{itemize}
    \column{0.46\textwidth}
      \textbf{Part 2 (90 min): Student Presentations}
      \begin{itemize}
        \item 3 presentations $\times$ 20 min + 10 min Q\&A
        \item Topics: RAG, agents, alignment, multimodal
        \item Peer evaluation
      \end{itemize}
  \end{columns}
  \vspace{0.4em}
  \textbf{Final reflection:} What should every practitioner know about foundation models and agentic AI?
\end{frame}

\section{Autonomous Research Agents}
\sectionframe{Autonomous Research Agents}{Towards AI that does science}

\begin{frame}{The AI Scientist Vision}
  \textbf{Lu et al.\ (2024)}~\cite{lu2024aiscientist}
  \vspace{0.4em}
  \begin{tcolorbox}[hitbox, title=The AI Scientist]
    A fully automated framework for scientific discovery:
    \begin{enumerate}
      \item \textbf{Ideation}: generate research ideas from literature.
      \item \textbf{Experiment execution}: write and run code, collect results.
      \item \textbf{Analysis}: interpret results, generate figures.
      \item \textbf{Paper writing}: produce a full NeurIPS-style paper.
      \item \textbf{Peer review}: self-evaluate using an LLM reviewer.
    \end{enumerate}
  \end{tcolorbox}
  \vspace{0.3em}
  \textbf{Results:} Papers produced on diffusion models, transformers, and RL at $\sim$\$15/paper.
  LLM reviewer scores $\sim$50\% on NeurIPS threshold.
\end{frame}

\begin{frame}{What AI Scientist Represents}
  \begin{columns}[T]
    \column{0.52\textwidth}
      \textbf{What it gets right:}
      \begin{itemize}
        \item End-to-end automation of repetitive scientific work.
        \item Rapid exploration of hyperparameter/architecture space.
        \item 24/7 operation at low marginal cost.
        \item Democratizes access to research capacity.
      \end{itemize}
    \column{0.44\textwidth}
      \textbf{What remains hard:}
      \begin{itemize}
        \item Genuinely \emph{novel} ideas (not incremental).
        \item Experimental rigor and statistical validity.
        \item Long-range planning across months.
        \item Physical experiments and wet lab integration.
        \item Ensuring results are honest, not cherry-picked.
      \end{itemize}
  \end{columns}
  \note{Ask students: would you trust a paper written by an AI scientist submitted to a journal? What safeguards would you want?}
\end{frame}

\section{Agentic Workflows in Production}
\sectionframe{Agentic Workflows in Production}{From research to deployed systems}

\begin{frame}{Production Agentic Patterns}
  \begin{center}
  \begin{tabular}{lll}
    \toprule
    \textbf{Pattern} & \textbf{Example} & \textbf{Key challenge} \\
    \midrule
    Coding assistant & GitHub Copilot Workspace & Code correctness \\
    Customer support & Zendesk AI, Sierra.ai & Escalation, hallucination \\
    Research assistant & Perplexity, Elicit & Citation accuracy \\
    DevOps automation & Incident.io, PagerDuty AI & Irreversible actions \\
    Data pipeline & Hex AI, Databricks AI & Schema drift \\
    Legal document & Harvey, Klarity & PII, jurisdiction \\
    \bottomrule
  \end{tabular}
  \end{center}
  \vspace{0.3em}
  \textbf{Common theme:} Every production deployment has added \emph{guardrails}, \emph{scoped tools},
  and \emph{human review} stages that research prototypes lack.
\end{frame}

\begin{frame}{Engineering for Production Agents}
  \begin{tcolorbox}[hitbox, title=Production Checklist]
    \begin{itemize}
      \item[$\square$] Observability: log every LLM call, tool call, and token count.
      \item[$\square$] Tracing: structured traces (e.g., LangSmith, Weights \& Biases Traces).
      \item[$\square$] Fallback policies: what happens when the LLM fails or times out?
      \item[$\square$] Cost controls: token budgets, rate limiting, model downgrade for cheap tasks.
      \item[$\square$] Human escalation: clear criteria for when to route to a human.
      \item[$\square$] Security review: red-team for prompt injection before launch.
      \item[$\square$] Evaluation suite: regression tests against a golden set of agent traces.
    \end{itemize}
  \end{tcolorbox}
\end{frame}

\section{Open Research Problems}
\sectionframe{Open Research Problems}{Where is the frontier?}

\begin{frame}{Open Problems in Foundation Models}
  \begin{columns}[T]
    \column{0.52\textwidth}
      \textbf{Architecture and training:}
      \begin{itemize}
        \item Efficient long-context models (beyond 1M tokens).
        \item Continual learning without catastrophic forgetting.
        \item Sample-efficient learning from few examples.
        \item Better calibration (model knows what it doesn't know).
      \end{itemize}
      \vspace{0.4em}
      \textbf{Reasoning:}
      \begin{itemize}
        \item Formal verification of LLM outputs.
        \item Neuro-symbolic integration.
        \item Compositional generalization.
      \end{itemize}
    \column{0.44\textwidth}
      \textbf{Agentic systems:}
      \begin{itemize}
        \item Robust prompt injection defense.
        \item Reliable multi-agent coordination.
        \item Long-horizon planning.
        \item Efficient memory management.
        \item Self-improving agents.
      \end{itemize}
      \vspace{0.3em}
      \textbf{Alignment:}
      \begin{itemize}
        \item Scalable oversight.
        \item Interpretability at scale.
        \item Value learning from sparse feedback.
      \end{itemize}
  \end{columns}
\end{frame}

\begin{frame}{Research Problems in Evaluation}
  \begin{tcolorbox}[hitbox, title=The Evaluation Crisis]
    \begin{itemize}
      \item Benchmarks are contaminated or gamed faster than new ones are created.
      \item LLM-as-judge has systematic biases (length, position, sycophancy).
      \item No ground truth for ``alignment'' — what does a well-aligned agent even look like?
      \item Agent evaluation is expensive: each eval run may cost hundreds of API calls.
    \end{itemize}
    \textbf{Needed:} Live, dynamic benchmarks; formal definitions of alignment;
    cheap-but-reliable automated evaluation.
  \end{tcolorbox}
\end{frame}

\section{Societal Implications}
\sectionframe{Societal Implications}{The broader picture}

\begin{frame}{Labor, Automation, and the Economy}
  \begin{columns}[T]
    \column{0.52\textwidth}
      \textbf{Jobs at risk (short term):}
      \begin{itemize}
        \item Content writing, translation, coding (partial).
        \item Customer support, data entry.
        \item Legal document review, financial analysis.
      \end{itemize}
      \vspace{0.4em}
      \textbf{Jobs created:}
      \begin{itemize}
        \item AI/ML engineers, data labelers.
        \item Prompt engineers, AI auditors.
        \item Human-in-the-loop specialists.
      \end{itemize}
    \column{0.44\textwidth}
      \textbf{Open questions:}
      \begin{itemize}
        \item Will AI create more jobs than it displaces?
        \item Who benefits from productivity gains?
        \item Will AI exacerbate inequality between AI-enabled and AI-excluded workers?
        \item How should educational institutions respond?
      \end{itemize}
  \end{columns}
\end{frame}

\begin{frame}{Power Concentration and Governance}
  \begin{tcolorbox}[alertbox, title=Concentration Risks]
    \begin{itemize}
      \item Frontier models require billions of dollars to train — accessible only to large labs.
      \item Cloud providers control inference infrastructure.
      \item Foundation models embed the values and biases of their creators at scale.
      \item National security implications: export controls on GPU chips.
    \end{itemize}
  \end{tcolorbox}
  \vspace{0.3em}
  \textbf{Governance approaches:}
  \begin{itemize}
    \item EU AI Act: risk-based regulation; transparency requirements.
    \item NIST AI RMF: voluntary framework for responsible AI.
    \item Open-source alternatives: Llama, Mistral, Falcon — distributes capability.
    \item International cooperation: AI safety institutes (UK, US, EU).
  \end{itemize}
\end{frame}

\begin{frame}{A Reflection: What This Course Has Covered}
  \begin{columns}[T]
    \column{0.50\textwidth}
      \textbf{Foundation:}
      \begin{itemize}
        \item Scaling laws and emergence (Wk 1)
        \item Transformer architecture (Wk 2)
        \item RLHF, DPO, alignment (Wk 3)
        \item Multimodal models (Wk 4)
      \end{itemize}
      \vspace{0.3em}
      \textbf{Retrieval:}
      \begin{itemize}
        \item Embeddings and semantic search (Wk 5)
        \item RAG and advanced retrieval (Wk 6)
      \end{itemize}
    \column{0.46\textwidth}
      \textbf{Agents:}
      \begin{itemize}
        \item Agentic AI foundations (Wk 7)
        \item Reasoning and planning (Wk 8)
        \item Memory systems (Wk 9)
        \item Multi-agent systems (Wk 10)
      \end{itemize}
      \vspace{0.3em}
      \textbf{Reliability:}
      \begin{itemize}
        \item Evaluation (Wk 11)
        \item Safety and security (Wk 12)
        \item Future directions (Wk 13)
      \end{itemize}
  \end{columns}
\end{frame}

\section{Student Presentations}
\sectionframe{Student Project Presentations}{Sharing your research}

\begin{frame}{Presentation Guidelines}
  \begin{tcolorbox}[hitbox, title=Format]
    \begin{itemize}
      \item \textbf{Duration:} 15–20 minutes presentation + 10 minutes Q\&A.
      \item \textbf{Slides:} 12–15 slides (title, problem, related work, method, results, discussion).
      \item \textbf{Demo:} short live demo or recorded video strongly encouraged.
      \item \textbf{Reproducibility:} share code and data (GitHub repo).
    \end{itemize}
  \end{tcolorbox}
\end{frame}

\begin{frame}{Presentation Evaluation Rubric}
  \begin{center}
  \begin{tabular}{lll}
    \toprule
    \textbf{Criterion} & \textbf{Weight} & \textbf{Description} \\
    \midrule
    Problem clarity & 15\% & Is the problem well-motivated? \\
    Technical depth & 25\% & Does the work show technical understanding? \\
    Experimental rigor & 20\% & Are results reliable and reproducible? \\
    Presentation quality & 20\% & Clear, well-paced, engaging \\
    Critical thinking & 20\% & Limitations, future work honestly discussed \\
    \bottomrule
  \end{tabular}
  \end{center}
  \vspace{0.3em}
  \textbf{Peer evaluation:} Each student rates 2 other presentations on a 1–5 scale.
  Peer scores contribute 20\% to the presentation grade.
\end{frame}

\section{Discussion \& Closing}
\begin{frame}{Final Discussion: Looking Forward}
  \begin{enumerate}
    \item What is the most important research contribution you would want to make to this field in the next 5 years?
    \item We covered many techniques that improve agent capability. Which improvement would have the most positive real-world impact? Which carries the most risk?
    \item Foundation models are trained by a handful of organizations but deployed by millions. What is the responsibility of a practitioner who deploys these models?
    \item If you had to redesign this course for 2030, what topics would you add? What would you remove?
  \end{enumerate}
\end{frame}

\begin{frame}{Course Summary and Takeaways}
  \begin{tcolorbox}[hitbox, title=Key Takeaways]
    \begin{enumerate}
      \item \textbf{Scale is not everything} — data quality, alignment, and architecture matter.
      \item \textbf{Retrieval augments knowledge} — RAG is more reliable than parametric memory alone.
      \item \textbf{Agents amplify both capability and risk} — every new tool is a new attack surface.
      \item \textbf{Evaluation is hard and critical} — don't trust benchmark numbers blindly.
      \item \textbf{Safety is not optional} — for deployed agents, security must be first-class.
      \item \textbf{The field evolves rapidly} — learn to read papers, not just apply frameworks.
    \end{enumerate}
  \end{tcolorbox}
  \vspace{0.3em}
  \textbf{Thank you for a great semester. Good luck with your projects and careers!}
\end{frame}

\begin{frame}[allowframebreaks]{Suggested Reading}
  \nocite{lu2024aiscientist, shen2023hugginggpt, bommasani2021foundation,
          anthropic2024claude3, wei2022emergent}
  \bibliography{../../references}
\end{frame}

\end{document}
